\section{Limits and colimits of presheaves}
\label{section-limits-colimits-PSh}

\noindent
Let $\mathcal{C}$ be a category.
Limits and colimits exist in the category
$\textit{PSh}(\mathcal{C})$. In addition, for any
$U \in \Ob(\mathcal{C})$ the functor
$$
\textit{PSh}(\mathcal{C})
\longrightarrow
\textit{Sets}, \quad
\mathcal{F}
\longmapsto
\mathcal{F}(U)
$$
commutes with limits and colimits. Perhaps the easiest way to prove
these statements is the following. Given a diagram
$
\mathcal{F} :
\mathcal{I}
\to
\textit{PSh}(\mathcal{C})
$
define presheaves
$$
\mathcal{F}_{\lim} :
U
\longmapsto
\lim_{i \in \mathcal{I}} \mathcal{F}_i(U)
\text{  and  }
\mathcal{F}_{\colim} :
U
\longmapsto
\colim_{i \in \mathcal{I}} \mathcal{F}_i(U)
$$
There are clearly projection maps $\mathcal{F}_{\lim} \to \mathcal{F}_i$
and canonical maps $\mathcal{F}_i \to \mathcal{F}_{\colim}$. These
maps satisfy the requirements of the maps of a limit (resp.\ colimit)
of Categories, Definition \ref{categories-definition-limit}
(resp.\ Categories, Definition \ref{categories-definition-colimit}).
Indeed, they clearly form a cone, resp. a cocone, over $\mathcal{F}$.
Furthermore, if $(\mathcal{G}, q_i : \mathcal{G} \to \mathcal{F}_i)$
is another
system (as in the definition of a limit), then we get for every
$U$ a system of maps $\mathcal{G}(U) \to \mathcal{F}_i(U)$
with suitable functoriality requirements. And thus a unique
map $\mathcal{G}(U) \to \mathcal{F}_{\lim}(U)$. It is easy
to verify these are compatible  as we vary $U$ and arise from
the desired map $\mathcal{G} \to \mathcal{F}_{\lim}$.
A similar argument works in the case of the colimit.





















